\slajdid{02-s002}
\begin{frame}[label=02-s002]{Šta znači „trenutno“ stanje ulaza?}
\fontsize{11}{12.5}\selectfont
% element: 02-s002:e01
Kombinacione mreže su digitalni sistemi kod kojih izlazi zavise samo od stanja ulaza. \alert{???}
\par\vspace{1.5mm}
Od kojih stanja ulaza? Jučerašnjih? Od onih pre 100 ms? Bivših, sadašnjih, budućih?\par
\textbf{Moramo dodati vremensku odrednicu!}
\par\vspace{1.5mm}
\Rezultat{Kombinacione mreže su digitalni sistemi kod kojih izlazi zavise samo od \alert{trenutnog} stanja ulaza.}
\par\vspace{1.5mm}
% element: 02-s002:e02
\textbf{Trenutnog?} Pojam koji često koristimo ne ulazeći dublje u njegovo značenje.\par
Ako su se ulazi promenili u nekom trenutku vremena, da li su se i izlazi \alert{trenutno} promenili?
\par\vspace{1.5mm}
Kada kažemo da se nešto trenutno dešava, obično podrazumevamo vremenske margine: u realnom fizičkom svetu promena se ne dešava u beskonačno kratkom vremenskom intervalu.
\end{frame}
